% Reference: https://en.cppreference.com/w/cpp/utility/pair.html
\subsection{pair、tuple 与结构化绑定}
\label{stl:pair-tuple}
\subsubsection*{原理}
\InlineCode{pair} 存两个值，\InlineCode{tuple} 可存多个不同类型的值。默认比较按字典序：先比较第一项，只有相等才比较下一项。因此数对可直接排序、放入有序容器或作为堆节点。
\subsubsection*{使用教程}
\InlineCode{p.first/second} 访问数对，\InlineCode{get<I>(t)} 访问元组的编译期第 $I$ 项（从 0 开始）。C++17 的 \InlineCode{auto [x,y]=p} 解包到副本；需修改原对象时用引用绑定。\InlineCode{tie(x,y)} 可同时赋值，也可用于多关键字比较。
\begin{lstlisting}[style=mtycpp,numbers=none]
using PII=pair<int,int>;
PII p={3,7};
auto [u,v]=p;
auto &[x,y]=p;
x+=1;
tuple<int,int,long long> t={1,2,9};
long long w=get<2>(t);
tie(u,v)=make_pair(5,6);
vector<PII> a={{2,3},{1,7},{2,1}};
sort(a.begin(),a.end());
\end{lstlisting}
若只按第一项比较，会把第一项相等的数对视为同一排序键；用于 \InlineCode{set} 时可能合并本想保留的元素。需要区分时把编号等次关键字也放入比较。
